\section{The execution model and recursive merge bases}
\label{sec:execution}

This section specifies the store underlying certified execution in
Definition~\ref{def:issuance}. A merge base is a property of the version graph;
its state is then used by the datatype's three-way merge.

\begin{definition}[Ranked store and ordinary execution]
\label{def:execution}
A configuration is $C=(S,h,\operatorname{par},\prec_C)$, where
$S(v)=(s_v,E_v)$ is a partial store, $h(r)$ is an active replica's allocated
head, and $\operatorname{par}(v)$ lists its parents. Version ranks are natural
numbers and each parent has smaller rank. Version $0$ stores
$(\sigma_0,\varnothing)$. Visibility endpoints belong to the events represented
at active heads. Visible predecessors of a head's events belong to that head,
distinct events have distinct timestamps, and
\[
 a\prec_C b\Longrightarrow a.t<b.t.
\]
Distinct events from the same replica are visibility-comparable. Reachable
stores additionally have transitive visibility and causally closed stored
histories. Write $v\preceq w$ for the reflexive transitive closure of parent
edges directed from parent to child.

All steps preserve the store except at the explicitly fresh version, preserve
other heads, and leave visibility unchanged unless stated otherwise:
\begin{enumerate}
\item \textbf{Fork.} For a fresh replica $r'$ and existing head $h(r)=v$,
allocate $v'>v$ with $S(v')=S(v)$, $\operatorname{par}(v')=[v]$, and
$h(r')=v'$. The source head is unchanged.
\item \textbf{Apply.} For $h(r)=v$, $S(v)=(s,E)$, and $e=(t,r,o)$,
allocate $v'>v$ with
\[
 S(v')=(u(s,e),E\cup\{e\}),\qquad
 \operatorname{par}(v')=[v],\qquad h(r)=v'.
\]
Replace visibility by $\prec_C\ \cup\ (E\times\{e\})$.
The timestamp is fresh against all events in the registry and every stored
history. The target configuration must satisfy the clock condition above,
so $p.t<t$ for every $p\in E$. Raw apply has no client guard; certified apply
adds $G(e,s)$ from Definition~\ref{def:issuance}.
\item \textbf{Ordinary merge.} For heads $v_1,v_2$ with stored states
$s_1,s_2$, choose an allocated greatest common ancestor $g$ as defined below.
Allocate $v'>v_1,v_2$ with
\[
 S(v')=(m(s_g,s_1,s_2),E_{v_1}\cup E_{v_2}),\qquad
 \operatorname{par}(v')=[v_1,v_2].
\]
Update the receiving replica's head to $v'$.
\item \textbf{Query.} Return $\operatorname{read}(s_{h(r)},q)$ without
changing the configuration.
\end{enumerate}
The initial store has just version $0$ and one replica pointing to it.
Executions are finite sequences of these steps. The certified relation also
retains the original-mint evidence described in Definition~\ref{def:issuance}.
\end{definition}
% Sources: Framework.Execution.Configuration, Step; Framework.Certificates.

\begin{definition}[Greatest common ancestor]
\label{def:gca}
A greatest common ancestor $g$ of $v_1,v_2$ satisfies
\[
 g\preceq v_1\land g\preceq v_2\land
 \forall w,\ (w\preceq v_1\land w\preceq v_2)\Longrightarrow w\preceq g.
\]
A maximal common ancestor need only have no strictly later common ancestor;
there may be several incomparable maximal common ancestors and no greatest one.
The store model carries the coherence property
$E_g=E_{v_1}\cap E_{v_2}$ whenever an allocated greatest common ancestor exists.
\end{definition}

\begin{theorem}[Intersection from event origins]
\label{thm:gca-intersection}
Suppose parents are allocated, histories grow along parent edges, and every
event has an allocated origin version containing it which is an ancestor of
\emph{every} version containing that event. Then for allocated $g,v_1,v_2$
with $g$ their greatest common ancestor,
\[
 E_g=E_{v_1}\cap E_{v_2}.
\]
\end{theorem}
The forward inclusion follows from history monotonicity. For a shared event,
its origin is common to both heads and therefore precedes $g$, proving the
reverse inclusion. The store-invariant proof supplies this justification for
the coherence field of Definition~\ref{def:gca}; it is not a statement about
arbitrary event labels attached to an arbitrary DAG.
% Source: Metatheory.StoreInvariant.gca_events_of_storeInv.

\begin{definition}[Recursive virtual merge base]
\label{def:virtual}
For a finite set of versions $T$ and a version $w$, let $M(T,w)$ be the list,
in increasing rank, of maximal elements of
\[
 \{x\mid x\preceq w\land\exists v\in T,\ x\preceq v\}.
\]
Write $s_v$ for the stored state at $v$, using $\sigma_0$ as the total-function
default when unallocated; correctness uses allocated supports. Define the
base $B$ and scratch fold $F$ by
\begin{align*}
 B(T,w)&=
 \begin{cases}
 \sigma_0,&M(T,w)=[],\\
 F(\{v\},s_v,\eta),&M(T,w)=v::\eta,
 \end{cases}\\
 F(T,s,[])&=s,\\
 F(T,s,v::\eta)&=F(T\cup\{v\},m(B(T,v),s,s_v),\eta).
\end{align*}
The virtual base of $v_1,v_2$ is $B(\{v_1\},v_2)$.
The widened merge step uses this state in place of $s_g$ in
Definition~\ref{def:execution}, without requiring a greatest common ancestor.
Only the final merged version is allocated. The scratch folds do not create
replicas, events, or intermediate stored versions.
\end{definition}
The recursion is well-founded by the size of the joint ancestor support,
lexicographically with the pending list length. The execution theorem must
establish correctness of these recursive scratch states as well as ordinary
stored versions; replacing $g$ by an arbitrary maximal common ancestor is not
this construction.
% Sources: Metatheory.VirtualMergeBase.vfoldAux, virtualBaseAux,
% virtualMergeBaseState; Framework.Execution.StepV.
